\documentclass[11pt,a4paper]{article}

\usepackage[utf8]{inputenc}
\usepackage[T1]{fontenc}
\usepackage[spanish]{babel}
\usepackage{amsmath,amssymb,amsfonts}
\usepackage{graphicx}
\usepackage{booktabs}
\usepackage{multirow}
\usepackage{hyperref}
\usepackage{xcolor}
\usepackage{listings}
\usepackage{algorithm}
\usepackage{algpseudocode}
\usepackage{geometry}
\usepackage{microtype}
\usepackage{natbib}
\usepackage{caption}
\usepackage{subcaption}
\usepackage{tikz}
\usetikzlibrary{arrows.meta,positioning,shapes.geometric}

\geometry{margin=2.5cm}

\hypersetup{
  colorlinks=true,
  linkcolor=blue!60!black,
  citecolor=green!40!black,
  urlcolor=blue!60!black,
}

\lstset{
  basicstyle=\ttfamily\footnotesize,
  breaklines=true,
  frame=single,
  backgroundcolor=\color{gray!8},
  keywordstyle=\color{blue!70!black},
  commentstyle=\color{gray},
  stringstyle=\color{orange!80!black},
  language=Python,
}

\title{%
  \textbf{Capibara Legal V2}: Un modelo de lenguaje compacto y especializado\\
  para el derecho español mediante tokenización BPE, razonamiento estructurado,\\
  destilación multi-semilla y recuperación jurídica específica de dominio%
}

\author{%
  Anachroni Co.\\
  \texttt{contacto@anachroni.co}
}

\date{Mayo 2026}

\begin{document}

\maketitle

\begin{abstract}
Presentamos \textbf{Capibara Legal V2}, un sistema de modelos de lenguaje compactos diseñado
para el dominio legal español. Partiendo de la arquitectura Transformer densa de V1
(vocab=512, seq=1024), V2 introduce ocho mejoras ortogonales que, combinadas, producen
un salto cualitativo sin requerir hardware acelerado: (1)~tokenizador BPE de 32\,000 tipos
entrenado sobre corpus jurídico español, cuadruplicando el contenido útil por ventana de
contexto; (2)~etiquetas Think-Anywhere para razonamiento interno estructurado; (3)~extensión
de seq\_len a 2048 tokens; (4)~wrapper OpenAI compatible con CrewAI para flujos de trabajo
multi-agente; (5)~fusión ponderada de adapters LoRA; (6)~sopa multi-semilla en destilación
con temperatura $T=4.0$; (7)~Infini-attention para contexto efectivamente ilimitado; y
(8)~modelo de embedding legal específico de dominio entrenado con la receta Arctic-Embed
y minería de negativos duros ajustable. El modelo Large V2 alcanza $\sim$515\,M parámetros
con vocab BPE y se despliega íntegramente en CPU ARM64 (Google Axion) sin GPU.
Describimos la arquitectura, el pipeline de entrenamiento y los objetivos de evaluación.
\end{abstract}

\tableofcontents
\newpage

%% ─────────────────────────────────────────────────────────────────────────────
\section{Introducción}
\label{sec:intro}
%% ─────────────────────────────────────────────────────────────────────────────

Los modelos de lenguaje grandes (LLMs) han demostrado capacidades notables en tareas
legales~\citep{nay2023large}, pero su despliegue en entornos profesionales reales
presenta obstáculos prácticos: coste de inferencia en GPU, ausencia de
especialización en terminología jurídica española, y limitaciones de contexto
insuficientes para documentos legales de varias páginas.

Los modelos de lenguaje pequeños (SLMs) especializados ofrecen una alternativa viable.
\citet{gurgurov2024adapting} demostraron que modelos BERT de 177\,M parámetros
adaptados a dominios específicos superan a LLaMA-3-8B en tareas de ese dominio.
\citet{belcak2025small} mostraron que un modelo Hymba de 1.5\,B parámetros supera
a modelos de 13\,B en instrucción-seguimiento con 3$\times$ el throughput.

Capibara Legal V2 extiende esta línea de trabajo con un enfoque específico al
derecho español, introduciendo ocho mejoras sobre V1 que actúan sobre el tokenizador,
la arquitectura, los datos de entrenamiento, el proceso de destilación y el sistema
de recuperación. Todas las mejoras son compatibles con despliegue en CPU sin GPU.

\paragraph{Contribuciones principales}
\begin{enumerate}
  \item Tokenizador BPE de 32\,000 tipos entrenado sobre corpus legal español,
        proporcionando cobertura 4$\times$ mayor por ventana de contexto.
  \item Esquema Think-Anywhere con tokens especiales reservados para razonamiento
        interno, con overhead del 4\% sobre el presupuesto de tokens.
  \item Protocolo de sopa multi-semilla en destilación con temperatura alta ($T=4.0$),
        mejorando la perplexity en 2--4\% respecto a la sopa intra-ejecución.
  \item Modelo de embedding legal \textit{capibara-embed-m} entrenado con
        estratificación por fuente y minería de negativos duros ajustable.
  \item Pipeline de mejora continua S1--S6 basado en logs de producción.
\end{enumerate}

%% ─────────────────────────────────────────────────────────────────────────────
\section{Arquitectura del sistema}
\label{sec:architecture}
%% ─────────────────────────────────────────────────────────────────────────────

\subsection{Familia de modelos}

V2 mantiene la topología cascada de V1, extendiendo los parámetros al introducir
el vocabulario BPE. La Tabla~\ref{tab:models} resume la familia de modelos.

\begin{table}[h]
\centering
\caption{Familia de modelos Capibara V2. Todos los modelos usan RoPE, pre-norm
y atención multi-cabeza con GQA opcional. El aumento de parámetros en V2 respecto a V1
se debe íntegramente a la tabla de embedding BPE (vocab=512 → 32\,000).}
\label{tab:models}
\begin{tabular}{lrrrrrl}
\toprule
\textbf{Nombre} & \textbf{$d$} & \textbf{$L$} & \textbf{$H$} & \textbf{seq} & \textbf{Parámetros} & \textbf{Rol} \\
\midrule
Cerebro   & 512  & 8  & 8  & 1024 & $\sim$42\,M  & Especulación (draft) \\
Medium    & 768  & 12 & 12 & 2048 & $\sim$138\,M & Verificación intermedia \\
Large     & 1280 & 24 & 20 & 2048 & $\sim$515\,M & Verificación final \\
\bottomrule
\end{tabular}
\end{table}

\subsection{Pipeline de inferencia especulativa}

V2 hereda el pipeline de decodificación especulativa de V1:

\begin{equation}
  \text{Cerebro} \xrightarrow{\text{draft } k \text{ tokens}} \text{Medium} \xrightarrow{\text{verify}} \text{Large}
\end{equation}

El modelo Cerebro genera $k=3$--5 tokens de forma especulativa; Medium verifica con
tasa de aceptación $\alpha \approx 87\%$; Large interviene solo cuando Medium rechaza.
Con think-tags en V2, los tokens \texttt{<think>}/\texttt{</think>} pasan por el mismo
pipeline sin modificación; se eliminan del output final antes de devolver al usuario.

\subsection{Adapters LoRA especializados}

El modelo Large sirve como base para 15 adapters LoRA entrenados por especialidad
jurídica y tipo de tarea (Tabla~\ref{tab:loras}).

\begin{table}[h]
\centering
\caption{Adapters LoRA V2. rank=16, $\alpha$=32 para todos excepto \textit{herramientas} y \textit{documentos}
(rank=8, $\alpha$=8, dropout=0.05, siguiendo~\citep{salesforce2025xlam}).}
\label{tab:loras}
\begin{tabular}{llll}
\toprule
\textbf{Categoría} & \textbf{Adapters} & \textbf{Configuración} & \textbf{Filtro} \\
\midrule
Legal por área  & penal, civil, laboral, constitucional, & rank=16, $\alpha$=32 & keywords \\
                & administrativo, mercantil & & \\
Tarea general   & resumen, instruccion, qa, extraccion, & rank=16, $\alpha$=32 & keywords \\
                & redaccion, dialogo, razonamiento, traduccion & & \\
Herramientas    & herramientas & rank=8, $\alpha$=8, dr=0.05 & JSON-parsability $\geq$99\% \\
Documentos      & documentos & rank=8, $\alpha$=8, dr=0.05 & estructura Markdown $\geq$0.4 \\
\bottomrule
\end{tabular}
\end{table}

%% ─────────────────────────────────────────────────────────────────────────────
\section{Mejora 1: Tokenizador BPE de 32k tipos}
\label{sec:bpe}
%% ─────────────────────────────────────────────────────────────────────────────

V1 usa tokenización byte-level (vocab=512): cada carácter español en UTF-8 ocupa
2 bytes, por lo que $\text{seq}=1024 \approx 200$--250 palabras. Con BPE entrenado
sobre corpus jurídico:

\begin{equation}
  1024 \text{ tokens}^{\text{BPE}} \approx 700\text{--}900 \text{ palabras} \quad
  (3.5\text{--}4\times \text{ más que byte-level})
\end{equation}

Términos como \textit{contencioso-administrativo} pasan de $\sim$25 tokens (byte-level)
a 2--3 (BPE legal). La tabla de embedding crece de $512 \times d$ a $32000 \times d$,
añadiendo 8--41\,M parámetros según el preset (Tabla~\ref{tab:emb}).

\begin{table}[h]
\centering
\caption{Impacto del tokenizador BPE en la tabla de embedding.}
\label{tab:emb}
\begin{tabular}{lrrl}
\toprule
\textbf{Preset} & \textbf{Emb V1} & \textbf{Emb V2 BPE} & \textbf{Params total V2} \\
\midrule
Cerebro  & 0.1\,M & 8.2\,M  & $\sim$42\,M  \\
Medium   & 0.4\,M & 24.6\,M & $\sim$138\,M \\
Large    & 0.7\,M & 40.9\,M & $\sim$515\,M \\
\bottomrule
\end{tabular}
\end{table}

Los tokens especiales \texttt{<think>} (ID 8) y \texttt{</think>} (ID 9) se reservan
en el vocabulario BPE, ocupando exactamente 1 token cada uno (overhead 4\% para bloques
de razonamiento de 50 tokens).

%% ─────────────────────────────────────────────────────────────────────────────
\section{Mejora 2: Think-Anywhere Tags}
\label{sec:think}
%% ─────────────────────────────────────────────────────────────────────────────

\subsection{Motivación}

El razonamiento en cadena (\textit{chain-of-thought}) mejora la precisión en tareas
jurídicas complejas~\citep{wei2022chain}. En V1, la sobrecarga de byte-level hace
inviable insertar bloques de razonamiento en el presupuesto de tokens. Con BPE y
tokens especiales reservados, el overhead es solo del 4\%.

\subsection{Arquitectura Think-Anywhere}

\begin{verbatim}
Consulta: "¿Puede un menor de 16 años firmar un contrato?"
[GENERACIÓN]
<think>
El artículo 1263 CC establece que los menores no emancipados
no pueden prestar consentimiento. La emancipación se produce
a los 16 años por varios supuestos (art. 314 CC).
</think>
Los menores de 16 años no emancipados no pueden firmar
contratos válidos según el artículo 1263 del Código Civil.
[OUTPUT VISIBLE]
Los menores de 16 años no emancipados...
\end{verbatim}

El bloque \texttt{<think>...\textbackslash</think>} se elimina del output mediante
expresión regular antes de devolver la respuesta al usuario.

\subsection{Datos de entrenamiento}

Siguiendo la \textit{Superficial Alignment Hypothesis}~\citep{zhou2023lima}, el
conocimiento jurídico proviene del pre-entrenamiento; el fine-tuning solo enseña
el formato de interacción. La Tabla~\ref{tab:think_data} muestra la distribución
objetivo de 1\,350 ejemplos curados.

\begin{table}[h]
\centering
\caption{Dataset de think traces curado (metodología LIMA).}
\label{tab:think_data}
\begin{tabular}{lrl}
\toprule
\textbf{Fuente} & \textbf{Ejemplos} & \textbf{Filtro} \\
\midrule
Stack Exchange Legal ES & 500 & $>$10 votos, respuesta aceptada, $>$200 palabras \\
Casos prácticos manuales & 200 & Revisión humana, formato IRAC \\
Think traces sintéticos (V1 teacher) & 600 & Perplexity $<$ umbral, verificado BOE \\
Diálogos multi-turno & 50 & $\geq$3 turnos, resolución clara \\
\midrule
\textbf{Total} & \textbf{1\,350} & \\
\bottomrule
\end{tabular}
\end{table}

Como datos de propósito general en español se utiliza UltraLink~\citep{ultralink2024}:
93\,K diálogos en español con capacidades matemáticas (32\,K), código (16\,K) y chat
(34\,K). El orden óptimo de SFT es inglés-reasoning primero, después español-legal,
maximizando la transferencia cross-lingual~\citep{ultralink2024}.

\subsection{Sampling adaptativo por especialidad}

Siguiendo~\citet{holtzman2020curious}, V2 implementa perfiles nucleus (top-p) por
adapter (Tabla~\ref{tab:sampling}).

\begin{table}[h]
\centering
\caption{Perfiles de sampling por especialidad.}
\label{tab:sampling}
\begin{tabular}{lcc}
\toprule
\textbf{Especialidad} & \textbf{Temperatura} & \textbf{top-p} \\
\midrule
extraccion, herramientas & 0.40 & 0.80 \\
razonamiento, qa         & 0.60 & 0.85 \\
resumen, legal (penal/civil/...) & 0.65 & 0.85 \\
documentos               & 0.70 & 0.88 \\
instruccion              & 0.75 & 0.90 \\
dialogo, redaccion       & 0.85 & 0.95 \\
\bottomrule
\end{tabular}
\end{table}

%% ─────────────────────────────────────────────────────────────────────────────
\section{Mejora 3: Extensión de contexto a seq\_len = 2048}
\label{sec:seqlen}
%% ─────────────────────────────────────────────────────────────────────────────

Con BPE, $\text{seq}=2048$ cubre 1400--1800 palabras, suficiente para:
un artículo legal completo, una sentencia corta o un contrato de 2--3 páginas.
La atención es $O(n^2)$: duplicar la secuencia cuadruplica la memoria de activaciones.
Se requiere gradient checkpointing obligatorio para Medium y Large.

\begin{table}[h]
\centering
\caption{Implicaciones de memoria y throughput al pasar a seq=2048.}
\label{tab:seq}
\begin{tabular}{lccccc}
\toprule
\textbf{Preset} & \textbf{seq V1} & \textbf{seq V2} & \textbf{$\Delta$ mem} & \textbf{Batch V2} & \textbf{Throughput est.} \\
\midrule
Cerebro & 512  & 1024 & $\times$4 & 16 & $\sim$2\,000 tok/s \\
Medium  & 1024 & 2048 & $\times$4 & 8  & $\sim$700 tok/s \\
Large   & 1024 & 2048 & $\times$4 & 4  & $\sim$150 tok/s \\
\bottomrule
\end{tabular}
\end{table}

%% ─────────────────────────────────────────────────────────────────────────────
\section{Mejora 4: Flujos multi-agente con CrewAI}
\label{sec:crewai}
%% ─────────────────────────────────────────────────────────────────────────────

V2 expone una API compatible OpenAI (\texttt{/v1/chat/completions}) mediante un
servidor de traducción de protocolo. Esto permite integrar Capibara con CrewAI
sin modificar el modelo:

\begin{center}
\texttt{CrewAI} $\xrightarrow{\text{OpenAI format}}$ \texttt{openai\_wrapper.py}
$\xrightarrow{\text{Capibara format}}$ \texttt{speculative\_inference.py}
\end{center}

Tres agentes configurados para flujos jurídicos: (1) \textit{Investigador Legal}
(búsqueda BOE/CENDOJ), (2) \textit{Analista Jurídico} (interpretación normativa),
(3) \textit{Redactor Legal} (generación de escritos).

%% ─────────────────────────────────────────────────────────────────────────────
\section{Mejora 5: Fusión ponderada de adapters LoRA}
\label{sec:lora_merge}
%% ─────────────────────────────────────────────────────────────────────────────

\subsection{Motivación}

Los 15 adapters LoRA de V1 son mutuamente excluyentes en cada inferencia.
LoRA merging combina varios adapters mediante media ponderada de las matrices
$B$ y $A$:

\begin{equation}
  W'_{\text{merged}} = W_{\text{base}} + \frac{\alpha}{r}
    \left(\sum_i w_i B_i A_i\right), \quad \sum_i w_i = 1
\end{equation}

\subsection{Adapters merged predefinidos}

\begin{table}[h]
\centering
\caption{Adapters merged para V2.}
\label{tab:merged}
\begin{tabular}{lll}
\toprule
\textbf{Adapter merged} & \textbf{Composición} & \textbf{Uso} \\
\midrule
\texttt{legal\_completo}     & 6 áreas legales ($w=1/6$)                     & Backup sin señal de routing \\
\texttt{razonamiento\_legal} & razonamiento (0.5) + penal/civil/laboral (0.5/3) & Análisis con CoT \\
\texttt{asistente\_legal}    & instruccion + dialogo + redaccion ($w=1/3$)   & Interfaz conversacional \\
\texttt{investigacion}       & herramientas + qa + extraccion ($w=1/3$)      & Búsqueda + extracción \\
\bottomrule
\end{tabular}
\end{table}

\subsection{Alineamiento de formato para herramientas}

Siguiendo~\citet{belcak2025small}, el adapter \textit{herramientas} se entrena
exclusivamente con ejemplos de JSON puro. Un filtro de parsabilidad JSON
(\textit{JSON parsability}) elimina ejemplos mixtos antes del fine-tuning:

\begin{equation}
  \text{JSON-parsability} = \frac{|\{y \in \mathcal{D} : \text{json.loads}(y) \neq \perp\}|}{|\mathcal{D}|}
\end{equation}

Objetivo: JSON-parsability $\geq 99\%$ tras fine-tuning LoRA con rank=8, $\alpha$=8,
dropout=0.05 y lr=$2\times10^{-5}$~\citep{salesforce2025xlam}.

\subsection{Alineamiento de formato para documentos}

El adapter \textit{documentos} genera escritos legales (contratos, demandas, actas,
poderes notariales, sentencias) en \textbf{Markdown estructurado} — el formato
nativo para exportación posterior a DOCX vía Pandoc, que es el estándar de los
despachos de abogados españoles.

Se define una métrica de estructura Markdown legal:

\begin{equation}
  s_{\text{doc}}(y) = \mathbb{1}[\text{headings}] \cdot 0.3
  + \mathbb{1}[\text{art.\,}N] \cdot 0.3
  + \mathbb{1}[\textbf{definición}] \cdot 0.2
  + \mathbb{1}[\text{lista num.}] \cdot 0.1
  + \mathbb{1}[|y| \geq 300] \cdot 0.1
\end{equation}

El filtro de entrenamiento exige $s_{\text{doc}} \geq 0.4$. Las convenciones de salida son:
secciones con \texttt{\#\#} / \texttt{\#\#\#}, artículos numerados (\texttt{Artículo 1.}),
definiciones en negrita, y ningún bloque de código fuente — solo texto legal fluido.

%% ─────────────────────────────────────────────────────────────────────────────
\section{Mejora 6: Sopa multi-semilla en destilación}
\label{sec:soup}
%% ─────────────────────────────────────────────────────────────────────────────

\subsection{Sopa de modelos y destilación}

\citet{wortsman2022model} demuestran que promediar pesos de varios modelos
fine-tuneados (\textit{model soups}) mejora la precisión sin coste de inferencia.
V1 ya aplica sopa intra-ejecución (últimos 3 checkpoints). V2 añade sopa inter-ejecución
y sopa de vectores de tarea.

\subsection{Fundamento teórico para la destilación}

La función objetivo de destilación con temperatura $T$:

\begin{equation}
  \mathcal{L}_{\text{distil}} = \alpha \cdot T^2 \cdot \text{KL}\!\left(
    \text{softmax}(z_T / T) \,\|\, \text{softmax}(z_S / T)
  \right) + (1 - \alpha) \cdot \mathcal{L}_{\text{CE}}
\end{equation}

con $T=4.0$ y $\alpha=0.7$. La temperatura alta produce un paisaje de loss más plano,
donde ejecuciones con distintas semillas convergen a mínimos locales en la misma cuenca.
El promedio de esos mínimos cae en una región de mayor generalización~\citep{wortsman2022model}.

\subsection{Jerarquía de 4 niveles}

\begin{equation}
\underbrace{
  \underbrace{
    \text{ckpt}_{a,t-2}, \text{ckpt}_{a,t-1}, \text{ckpt}_{a,t}
  }_{\text{intra-run}} \!\!\!\!\to \text{soup}_{a}
}_{\text{nivel 1}} \quad
\underbrace{
  \text{soup}_{42}, \text{soup}_{123}, \text{soup}_{777} \to \text{soup}_{\text{multi-seed}}
}_{\text{nivel 2}} \quad
\underbrace{
  \text{soup}_{\text{UltraLink}}, \text{soup}_{\text{legal}} \to \text{model}_{\text{final}}
}_{\text{nivel 3, task-vector}}
\end{equation}

El nivel 3 implementa fusión de vectores de tarea~\citep{ilharco2023editing}:

\begin{equation}
  \theta_{\text{final}} = \theta_{\text{base}} + \sum_i \lambda_i \cdot (\theta_{\text{soup}_i} - \theta_{\text{base}})
\end{equation}

con $\lambda_{\text{UltraLink}}=0.25$ y $\lambda_{\text{legal}}=0.75$ para priorizar
la especialización jurídica sobre el español general.

\subsection{Ganancia esperada}

\begin{table}[h]
\centering
\caption{Coste vs ganancia en distintas variantes de sopa para destilación.}
\label{tab:soup}
\begin{tabular}{lcc}
\toprule
\textbf{Variante} & \textbf{Coste} & \textbf{Ganancia en perplexity} \\
\midrule
V1: intra-run (3 ckpts)  & 1$\times$ destilación & baseline \\
V2: intra-run $\times$ 3 semillas & 3$\times$ destilación & $-$2--4\% \\
V2 opt.: 5 semillas & 5$\times$ destilación & $-$3--5\% (rendimientos decrecientes) \\
\bottomrule
\end{tabular}
\end{table}

%% ─────────────────────────────────────────────────────────────────────────────
\section{Mejora 7: Contexto extendido con Infini-attention}
\label{sec:infini}
%% ─────────────────────────────────────────────────────────────────────────────

\subsection{Infini-attention}

\citet{munkhdalai2024leave} proponen añadir una memoria compresiva asociativa
$M \in \mathbb{R}^{d_k \times d_v}$ por cabeza de atención. La actualización sigue
la regla Delta:

\begin{align}
  A_{\text{mem}} &= \frac{\sigma(Q) M_{s-1}}{\sigma(Q) z_{s-1}} \\
  M_s &= M_{s-1} + \sigma(K)^\top \left(V - \frac{\sigma(K) M_{s-1}}{\sigma(K) z_{s-1}}\right) \\
  A &= \text{sigmoid}(\beta) \cdot A_{\text{mem}} + (1 - \text{sigmoid}(\beta)) \cdot A_{\text{dot}}
\end{align}

donde $\sigma = \text{ELU}+1$, $z$ es el normalizador y $\beta$ es un escalar
aprendido por cabeza. La memoria $M$ tiene tamaño fijo independiente de $N$, frente
al KV-cache de los Transformers estándar que crece linealmente.

\subsection{Comparación de footprint de memoria}

\begin{table}[h]
\centering
\caption{Memoria de contexto en distintas arquitecturas~\citep{munkhdalai2024leave}.}
\label{tab:infini}
\begin{tabular}{lll}
\toprule
\textbf{Arquitectura} & \textbf{Memoria de contexto} & \textbf{Contexto efectivo} \\
\midrule
Transformer-XL & crece con $N \times l$ & Solo último segmento \\
Memorizing Transformers (65K KV) & 183\,M parámetros & 65K tokens \\
Infini-Transformer & 1.6\,M (114$\times$ menor) & Ilimitado (acumulativo) \\
\bottomrule
\end{tabular}
\end{table}

\subsection{Position Interpolation como paso previo}

Antes del continual pre-training de Infini-attention ($\sim$3 días), se aplica
Position Interpolation~\citep{chen2023extending} para extender $\text{seq}=2048$ a
$\text{seq}=16384$ con solo 1000 pasos ($\sim$2 horas):

\begin{equation}
  f'(x, m) = f\!\left(x,\; m \cdot \frac{L}{L'}\right), \quad L=2048,\; L'=16384
\end{equation}

Resultados sobre LLaMA 7B--65B: degradación $<2\%$ en benchmarks originales,
passkey retrieval al 100\% tras 200 pasos~\citep{chen2023extending}.

%% ─────────────────────────────────────────────────────────────────────────────
\section{Mejora 8: Embedding legal específico de dominio}
\label{sec:embed}
%% ─────────────────────────────────────────────────────────────────────────────

\subsection{Limitaciones de embeddings genéricos para RAG jurídico}

Los modelos de embedding multilingüe genéricos (multilingual-e5, LaBSE) no capturan:
terminología legal específica (\textit{interdicto}, \textit{apremio}), citas
jurídicas (STS 234/2021 vs SAP Madrid 44/2021), ni la jerarquía de fuentes
(constitución $\succ$ ley orgánica $\succ$ ley ordinaria $\succ$ reglamento).

\subsection{Receta Arctic-Embed para el dominio legal}

\citet{merrick2024arctic} identifican tres innovaciones sobre métodos anteriores:

\paragraph{1. Source stratification.} Cada mini-batch contiene ejemplos de una sola
fuente ($+3.2$ puntos nDCG@10 vs sin estratificación). Las fuentes legales son
naturalmente estratificables: BOE, CENDOJ, TJUE, legislación LATAM, académico,
sintético.

\paragraph{2. Minería de negativos duros ajustable (\textit{tunable hard negatives}).}

\begin{equation}
  \mathcal{N}^- = \{d \in \mathcal{C} : R_{\min} \leq \text{sim}(q, d) \leq R_{\max}\}
\end{equation}

con $R_{\min}=0.1$ y $R_{\max}=0.7$. En derecho, $R_{\max}=0.7$ es crítico:
artículos del mismo código con estructura similar pero contenido distinto
deben ser negativos, no positivos.

\paragraph{3. Generación sintética de queries fundamentada en negativos.}

Se usa el LLM V2 para generar consultas que recuperen el documento positivo pero
\textit{no} los documentos negativos duros, produciendo consultas con mayor
discriminación que la generación de queries sobre documentos aislados.

\subsection{Resultados esperados}

\begin{table}[h]
\centering
\caption{Ganancia esperada en retrieval legal con \textit{capibara-embed-m}.}
\label{tab:retrieval}
\begin{tabular}{lcc}
\toprule
\textbf{Tarea} & \textbf{Antes (e5 genérico)} & \textbf{Esperado (capibara-embed)} \\
\midrule
Recuperar artículo por descripción & $\sim$60\% R@10 & $\sim$80\% R@10 \\
Recuperar sentencia por holding    & $\sim$50\% R@10 & $\sim$75\% R@10 \\
Discriminar artículos similares    & $\sim$40\% nDCG@10 & $\sim$65\% nDCG@10 \\
Recuperar en LATAM (ES-MX, ES-AR)  & $\sim$45\% R@10 & $\sim$70\% R@10 \\
\bottomrule
\end{tabular}
\end{table}

%% ─────────────────────────────────────────────────────────────────────────────
\section{Pipeline de mejora continua S1--S6}
\label{sec:s1s6}
%% ─────────────────────────────────────────────────────────────────────────────

Siguiendo~\citet{belcak2025small}, el bucle S1--S6 convierte logs de producción
en datos de SFT para mejorar continuamente los adapters (Tabla~\ref{tab:s1s6}).
El umbral de re-entrenamiento es 200 ejemplos curados por adapter, con frecuencia
máxima de 1 re-entrenamiento cada 2 semanas.

\begin{table}[h]
\centering
\caption{Mapeo del algoritmo S1--S6 a la infraestructura de Capibara.}
\label{tab:s1s6}
\begin{tabular}{clll}
\toprule
\textbf{Paso} & \textbf{Descripción} & \textbf{Implementación} \\
\midrule
S1 & Registrar cada llamada al LLM & \texttt{log\_inference()} en servidor HTTP \\
S2 & Curar interacciones & \texttt{scripts/curate\_instruction\_data.py} \\
S3 & Clasificar por tarea & \texttt{\_route\_specialty()} en speculative\_inference \\
S4 & Seleccionar SLM apropiado & Slim200M + LoRA adapter por especialidad \\
S5 & Fine-tuning con LoRA & \texttt{scripts/lora\_finetune.py} \\
S6 & Iterar & Retorna a S1 \\
\bottomrule
\end{tabular}
\end{table}

%% ─────────────────────────────────────────────────────────────────────────────
\section{Plan de entrenamiento y despliegue}
\label{sec:training}
%% ─────────────────────────────────────────────────────────────────────────────

\subsection{Hardware y estimaciones}

Todo el entrenamiento se ejecuta en CPU ARM64 (Google Axion c4a-standard-32,
32 vCPU Neoverse V2, 128\,GB RAM) sin GPU. La Tabla~\ref{tab:timeline} resume
las estimaciones de tiempo.

\begin{table}[h]
\centering
\caption{Estimación del pipeline de entrenamiento V2 en Axion (32 threads).}
\label{tab:timeline}
\begin{tabular}{llrc}
\toprule
\textbf{Fase} & \textbf{Tarea} & \textbf{Steps} & \textbf{Tiempo est.} \\
\midrule
1 & Tokenizador BPE  & -- & 1--3 h \\
2 & Cerebro V2 (42M) & 15\,000 & 21 h \\
2 & Medium V2 (138M) & 10\,000 & 40 h \\
3 & Large V2 fase 1 (515M) & 35\,000 & 14 d \\
3 & Large V2 DAPT legal & 10\,000 & 4 d \\
4 & Position Interpolation $\times$ 1\,000 & 1\,000 & 2 h \\
4 & Destilación $\times$ 3 seeds $\times$ 3 targets & 30\,000 & 15 d \\
5 & LoRA $\times$ 15 adapters & 30\,000 & 24 h \\
5 & LoRA merging + benchmarks & -- & 8 h \\
\midrule
  & \textbf{Total} & & $\sim$\textbf{37 días} \\
\bottomrule
\end{tabular}
\end{table}

\subsection{Velocidad de inferencia teórica}

Con bf16 en 32 CPUs ARM Neoverse V2 ($\sim$150\,GB/s bandwidth efectivo), la
inferencia autoregresiva a batch=1 está limitada por bandwidth de memoria:

\begin{equation}
  \text{tok/s} \approx \frac{\text{bandwidth}}{\text{tamaño modelo bf16}}
\end{equation}

\begin{table}[h]
\centering
\caption{Velocidad de inferencia teórica en Axion (batch=1, autoregresivo).}
\label{tab:speed}
\begin{tabular}{lccc}
\toprule
\textbf{Modelo} & \textbf{Tamaño bf16} & \textbf{tok/s} & \textbf{Con spec. decoding} \\
\midrule
Cerebro (42M)  & 84\,MB   & $\sim$1\,785 tok/s & --               \\
Medium (138M)  & 276\,MB  & $\sim$543 tok/s    & --               \\
Large (515M)   & 1.03\,GB & $\sim$146 tok/s    & $\sim$360 tok/s  \\
\bottomrule
\end{tabular}
\end{table}

%% ─────────────────────────────────────────────────────────────────────────────
\section{Evaluación}
\label{sec:eval}
%% ─────────────────────────────────────────────────────────────────────────────

\begin{table}[h]
\centering
\caption{Objetivos de evaluación V1 $\to$ V2.}
\label{tab:benchmarks}
\begin{tabular}{lcc}
\toprule
\textbf{Métrica} & \textbf{V1 baseline} & \textbf{V2 objetivo} \\
\midrule
Perplexity legal (test set) & -- & $<$ V1 \\
Tokens/s Large+LoRA & $\sim$350 tok/s & $\geq$350 tok/s \\
Tasa de aceptación Medium spec. & $\sim$87\% & $\geq$87\% \\
BLEU summarization & -- & $>$ V1 \\
Exact match QA legal & -- & $>$ V1 \\
Palabras efectivas en contexto & $\sim$220 (seq 1024 byte-level) & $\sim$1500 (seq 2048 BPE) \\
Ganancia sopa multi-semilla & -- & $-$2--4\% perplexity \\
JSON-parsability herramientas & $\sim$0--7\% (zero-shot) & $\geq$99\% (post-LoRA) \\
$s_{\text{doc}}$ documentos (estructura Markdown) & $\sim$0.1 (zero-shot) & $\geq$0.75 (post-LoRA) \\
Passkey retrieval @32K tokens & 0\% & $\sim$100\% (Infini) \\
R@10 retrieval legal & $\sim$60\% & $\sim$80\% (capibara-embed) \\
\bottomrule
\end{tabular}
\end{table}

%% ─────────────────────────────────────────────────────────────────────────────
\section{Trabajo relacionado}
\label{sec:related}
%% ─────────────────────────────────────────────────────────────────────────────

\paragraph{LLMs para derecho.} \citet{nay2023large} revisan aplicaciones de LLMs
en tareas jurídicas. \citet{zheng2023does} estudian GPT-4 como asistente legal.
Nuestro trabajo se diferencia por el enfoque en despliegue CPU-only y el dominio
específicamente español.

\paragraph{Modelos compactos especializados.} \citet{gurgurov2024adapting} demuestran
superioridad de modelos especializados pequeños sobre modelos grandes generalistas
en dominios específicos. \citet{belcak2025small} extienden este resultado a pipelines
agénticos multi-SLM.

\paragraph{Destilación y sopa de modelos.} \citet{hinton2015distilling} introduce
la destilación con temperatura. \citet{wortsman2022model} proponen model soups.
\citet{ilharco2023editing} introducen task vectors para fusión de especializaciones.

\paragraph{Contexto largo.} \citet{munkhdalai2024leave} proponen Infini-attention.
\citet{chen2023extending} introducen Position Interpolation para extender contexto.
\citet{bertsch2024unlimiformer} y \citet{gu2023mamba} ofrecen alternativas SSM.

\paragraph{Modelos de embedding.} \citet{merrick2024arctic} introducen la receta
Arctic-Embed con estratificación por fuente y minería de negativos duros.

%% ─────────────────────────────────────────────────────────────────────────────
\section{Conclusiones}
\label{sec:conclusion}
%% ─────────────────────────────────────────────────────────────────────────────

Capibara Legal V2 demuestra que es posible construir un sistema de asistencia legal
en español de alta calidad íntegramente sobre CPU mediante la combinación de:
(1) tokenización BPE específica de dominio, (2) razonamiento estructurado con baja
sobrecarga, (3) destilación multi-semilla con sopa de vectores de tarea, y (4) retrieval
jurídico específico de dominio. La arquitectura cascada con decodificación especulativa
mantiene latencias competitivas ($\geq$360\,tok/s efectivos para Large) sin GPU.

El sistema sirve como base para V3, que introducirá arquitecturas híbridas
Mamba-Attention con Mixture-of-Experts y un corpus iberoamericano de $\sim$16\,B
tokens.

%% ─────────────────────────────────────────────────────────────────────────────
\bibliographystyle{plainnat}
\begin{thebibliography}{99}

\bibitem[Belcak et al.(2025)]{belcak2025small}
Belcak, P., Wattenhofer, R., et al. (2025).
\newblock Small Language Models are the Future of Agentic AI.
\newblock \textit{arXiv:2506.02153v2}.

\bibitem[Chen et al.(2023)]{chen2023extending}
Chen, S., Wong, S., Chen, L., and Tian, Y. (2023).
\newblock Extending Context Window of Large Language Models via Position Interpolation.
\newblock \textit{arXiv:2306.15595v2}, Meta AI.

\bibitem[Gurgurov et al.(2024)]{gurgurov2024adapting}
Gurgurov, D., Hartmann, M., and Ostling, R. (2024).
\newblock Adapting LLMs to Low-Resource Languages.
\newblock \textit{arXiv}.

\bibitem[Hinton et al.(2015)]{hinton2015distilling}
Hinton, G., Vinyals, O., and Dean, J. (2015).
\newblock Distilling the Knowledge in a Neural Network.
\newblock \textit{arXiv:1503.02531}.

\bibitem[Holtzman et al.(2020)]{holtzman2020curious}
Holtzman, A., Buys, J., Du, L., Forbes, M., and Choi, Y. (2020).
\newblock The Curious Case of Neural Text Degeneration.
\newblock \textit{ICLR 2020}.

\bibitem[Ilharco et al.(2023)]{ilharco2023editing}
Ilharco, G., Ribeiro, M. T., Wortsman, M., Gururangan, S., Schmidt, L., Hajishirzi, H.,
and Farhadi, A. (2023).
\newblock Editing Models with Task Arithmetic.
\newblock \textit{ICLR 2023}.

\bibitem[Merrick et al.(2024)]{merrick2024arctic}
Merrick, L., et al. (2024).
\newblock Arctic-Embed: Scalable, Efficient, and Accurate Text Embedding Models.
\newblock \textit{arXiv:2405.05374v1}, Snowflake Inc.

\bibitem[Munkhdalai et al.(2024)]{munkhdalai2024leave}
Munkhdalai, T., Faruqui, M., and Gopal, S. (2024).
\newblock Leave No Context Behind: Efficient Infinite Context Transformers with Infini-attention.
\newblock \textit{arXiv:2404.07143}, Google.

\bibitem[Nay et al.(2023)]{nay2023large}
Nay, J. J., et al. (2023).
\newblock Large Language Models as Tax Attorneys: A Case Study in Legal Capabilities.
\newblock \textit{arXiv:2306.07032}.

\bibitem[Salesforce(2025)]{salesforce2025xlam}
Salesforce Research (2025).
\newblock XLAM: A Family of Large Action Models to Empower AI Agent Systems.
\newblock \textit{arXiv:2504.19277v1}.

\bibitem[UltraLink(2024)]{ultralink2024}
Zhang, C., et al. (2024).
\newblock UltraLink: An Open-Source Knowledge-Enhanced Multilingual Supervised Fine-tuning Dataset.
\newblock \textit{arXiv:2402.04588}, Tsinghua University.

\bibitem[Wei et al.(2022)]{wei2022chain}
Wei, J., et al. (2022).
\newblock Chain-of-Thought Prompting Elicits Reasoning in Large Language Models.
\newblock \textit{NeurIPS 2022}.

\bibitem[Wortsman et al.(2022)]{wortsman2022model}
Wortsman, M., et al. (2022).
\newblock Model Soups: Averaging Weights of Multiple Fine-tuned Models Improves Accuracy
without Increasing Inference Time.
\newblock \textit{ICML 2022}.

\bibitem[Zhou et al.(2023)]{zhou2023lima}
Zhou, C., et al. (2023).
\newblock LIMA: Less Is More for Alignment.
\newblock \textit{NeurIPS 2023}, \textit{arXiv:2305.11206v1}.

\end{thebibliography}

\end{document}
